小 batch 可能没有足够的查询来占满 GPU。把一个查询的候选分给多个 block 能提供更多并行工作，但这些 block 必须共同保留全局最佳候选。因此，各 block 需要完整短名单的容量，再归并各列表 \cite{johnson_gpu}。这里的精确性使用共同的分数与位置顺序；只按距离截断时，并列候选可能不同。

\subsection{候选分区与确定性排序}

对一个查询，按 probe 顺序将选定 IVF 列表串接为 $n_c$ 个候选出现实例。假设每个被访问行只出现一次，且索引保持固定。对 $\gscan\in\{1,2,4,8\}$，block $(q,j)$ 扫描
\begin{equation}
 U_j=\left[\left\lfloor\frac{j n_c}{\gscan}\right\rfloor,
              \left\lfloor\frac{(j+1)n_c}{\gscan}\right\rfloor\right),
 \qquad j=0,\ldots,\gscan-1.
 \label{eq:5-partition}
\end{equation}
区间表示串接序列中的候选，而非数据库 ID。各区间大小最多相差一项，长列表可能跨越边界。这比按相同列表个数分配更能平衡候选数量。probe 前缀 offsets 定位每线程的首个列表，随后以游标推进至后续列表；本阶段只读取主码。

为了使局部和归并选择一致，相同分数按物理位置排序：
\begin{equation}
 K_q(x)=\bigl(\primarydist(q,x),p(x)\bigr),
 \qquad
 (a,p)\prec(b,r)\iff a<b\ \text{or}\ (a=b\ \text{and}\ p<r),
 \label{eq:5-primary-key}
\end{equation}
位置 $p(x)$ 在固定索引中唯一，也不依赖原子追加顺序。对非 NaN 分数，这个完整键可以确定所有候选的先后。最终结果的并列则使用数据库 ID。只要各路径计算分数和排序规则相同，改变分区就不改变选中的键。

\subsection{精确局部保留与全局覆盖}

以 $\operatorname{Top}_c(V)$ 表示前 $\min(c,|V|)$ 个键，其中 $c=\max\{k,\lceil\alpha k\rceil\}$。所有全局优胜者都可能落在同一组，因而每组只给 $c/\gscan$ 配额并不安全；每组须有保留至多 $c$ 项的容量。下面的缓冲保证说明：压缩保留仍符合共享界的局部优胜者。共享界 $\beta_t$ 在 \cref{sec:5-5} 介绍；不共享时，它是排在一切有效键之后的占位键，称为哨兵。

\begin{resultblock}
\begin{lemma}[局部缓冲保留]
\label{lem:5-buffer}
令 $V_t$ 为一个 block 扫描到 chunk $t$ 为止已评分的候选，$B_t$ 为该时刻任意压缩或过滤后的有效缓冲。假设完整键固定、追加容量充足且 top-$c$ 压缩精确。令 $\beta_t$ 为迄今观察到的完整共享阈值的最小值，包括初始顶部哨兵。则
\begin{equation}
 \operatorname{Top}_c(V_t)\cap\{x:K_q(x)\preceq\beta_t\}
 \subseteq B_t.
 \label{eq:5-buffer-invariant}
\end{equation}
不共享时，阈值为哨兵，缓冲包含整个前缀 top-$c$。
\end{lemma}
\begin{proof}
空前缀显然成立。令
\[
 E_t=\operatorname{Top}_c(V_t)\cap\{x:K_q(x)\preceq\beta_t\}.
\]
将 $V_t$ 扩展为 $V_{t+1}$ 时，$\beta_{t+1}\preceq\beta_t$，且
\[
 \operatorname{Top}_c(V_{t+1})\cap V_t\subseteq\operatorname{Top}_c(V_t).
\]
因此，
\[
 x\in E_{t+1}\cap V_t\Rightarrow x\in E_t\subseteq B_t.
\]
局部拒绝的新键已有 $c$ 个前驱，故排名 $>c$。容量可容纳其余符合准入的到达项。对任意符合条件的优胜者 $x$，
\[
 \#\{y\in B:K_q(y)\prec K_q(x)\}
 \leq\#\{y\in V_{t+1}:K_q(y)\prec K_q(x)\}<c,
\]
所以精确压缩保留 $x$。过滤移除 $K>\beta_{\rm seen}$；运行最小值覆盖陈旧读取。上述步骤通过归纳保持不变量。
\end{proof}
\end{resultblock}
一旦保留列表覆盖全部全局优胜者，就能从中恢复精确的 top-$c$。下一命题还允许每次归并只留下前 $c$ 项。

\begin{resultblock}
\begin{proposition}[覆盖与截断归并]
\label{prop:5-local-coverage}
在共同全序下固定键，令 $U$ 为已访问空间，$V$ 为任意保留子集。若 $\operatorname{Top}_c(U)\subseteq V\subseteq U$，则 $\operatorname{Top}_c(V)=\operatorname{Top}_c(U)$。成对有序归并并截断为 $c$ 可保持该结果。特别地，对不相交分组，
\begin{equation}
 \operatorname{Top}_c\!\left(\bigcup_j U_j\right)
 =\operatorname{Top}_c\!\left(\bigcup_j
                    \operatorname{Top}_c(U_j)\right).
 \label{eq:5-topc-union}
\end{equation}
该陈述也包含候选空间小于 $c$ 的情形。
\end{proposition}
\begin{proof}
每个全局优胜者都在 $V$ 内，而 $V$ 的其他成员均排在全局优胜者之后。因此，从 $V$ 选择可找回它们。对独立分组这一特例，设 $x$ 属于全局 top-$c$，且位于组 $U_j$。全局上排在 $x$ 前面的候选少于 $c$，所以在 $U_j$ 内排在它前面的也少于 $c$。于是 $x\in\operatorname{Top}_c(U_j)$，局部结果的并集包含所有全局优胜者。取并集前 $c$ 项，即按相同顺序精确找回全局优胜者。若整个空间少于 $c$ 个成员，则全部保留。

归并时，两个有序列表的前 $c$ 项就是并集的 top-$c$。被丢弃候选在该并集中已有至少 $c$ 个前驱，无法进入任何更大并集的 top-$c$。对归并树归纳即可得到结论。
\end{proof}
\end{resultblock}
\par\medskip\noindent\begin{minipage}{\linewidth}
\captionsetup{hypcap=false}
\centering
% Schematic only: illustrative keys, not experimental observations.
\begin{tikzpicture}[
  font=\small\linespread{1}\selectfont,
  box/.style={draw=ThesisNavy!65,rounded corners=2pt,align=center,
    inner sep=5pt,text width=2.65cm,minimum height=2.05cm},
  merged/.style={box,text width=5.55cm,minimum height=1.05cm,
    fill=ThesisNavy!5},
  arr/.style={-{Latex[length=2mm]},draw=ThesisNavy!80,thick}]
\node[box] (g0) at (0,0) {$\operatorname{Top}_3(U_0)$\\[2pt]
  $\mathbf{(1,8)}$\\$\mathbf{(2,6)}$\\$(2,9)$};
\node[box] (g1) at (3.5,0) {$\operatorname{Top}_3(U_1)$\\[2pt]
  $(4,1)$\\$(5,4)$\\\textit{length 2}};
\node[box] (g2) at (7,0) {$\operatorname{Top}_3(U_2)$\\[2pt]
  $\varnothing$\\\textit{length 0}};
\node[box] (g3) at (10.5,0) {$\operatorname{Top}_3(U_3)$\\[2pt]
  $\mathbf{(2,3)}$\\$(7,2)$\\$(8,5)$};
\node[merged] (m0) at (1.75,-2.5) {Merge and retain first 3\\[2pt]
  $\mathbf{(1,8)},\ \mathbf{(2,6)},\ (2,9)$};
\node[merged] (m1) at (8.75,-2.5) {Merge and retain first 3\\[2pt]
  $\mathbf{(2,3)},\ (7,2),\ (8,5)$};
\node[merged,text width=6.8cm,fill=ThesisNavy!12] (out) at (5.25,-4.3)
  {Global primary top-3\\[2pt]
   $\mathbf{(1,8)},\ \mathbf{(2,3)},\ \mathbf{(2,6)}$};
\draw[arr] (g0.south) -- (m0.north west);
\draw[arr] (g1.south) -- (m0.north east);
\draw[arr] (g2.south) -- (m1.north west);
\draw[arr] (g3.south) -- (m1.north east);
\draw[arr] (m0.south) -- (out.north west);
\draw[arr] (m1.south) -- (out.north east);
\end{tikzpicture}

\captionof{figure}[局部覆盖与截断成对归并]{局部覆盖与截断成对归并，示例取 $c=3$、$\gscan=4$。条目是示意的 $(\primarydist,p)$ 键，粗体为全局优胜者。短组／空组贡献有效前缀。归并按相同顺序最多保留 $c$ 项。}
\label{fig:5-global-selection}
\end{minipage}\par\medskip
每个 block 在共享内存中保存容量为 $C_{\rm buf}$ 的键对缓冲，该容量为二的幂。对 $T_{\rm blk}$ 个扫描线程，预留 $C_{\rm buf}\geq c+T_{\rm blk}$ 项，容纳保留键及下一个 chunk。bitonic 压缩排序缓冲并最多保留 $c$ 个有效键。有 $c$ 个键后，最差保留键成为局部阈值 $\lambda_j$；此前阈值为哨兵。新候选只有在 $K\prec\lambda_j$ 时进入。这 $c$ 个已保留键使候选继续到达时阈值仍然有效。

压缩之后，一个 chunk 每线程最多添加一个键，因此
\[
 |B|+|\Delta V|\leq c+T_{\rm blk}\leq C_{\rm buf}.
\]
这给出 \cref{lem:5-buffer} 的追加容量前提。结合 \cref{prop:5-shared-bound,prop:5-local-coverage}，便将缓冲实现连接到全局覆盖。写出全部主码分数需要 $O(n_c)$ 的存储和写入；扫描中直接保留候选则输出 $O(\gscan c)$ 个键对。寄存器选择也是另一种方案 \cite{johnson_gpu}。

\subsection{归并实现、存储与边界情况}

对 $\gscan>1$，一个 256 线程 block 为每查询执行 $\log_2\gscan$ 轮归并。长度为 $a,b$ 的一对列表通过 Merge Path 切片输出 $\min(c,a+b)$ 个键 \cite{gpu_merge_path}。比较遵循 \cref{eq:5-primary-key}，轮间使用 block 屏障。\Cref{fig:5-global-selection} 展示截断为何安全。除划分归并工作的搜索和同步开销外，内部节点最多输出 $c(\gscan-1)$ 个键对。

两块缓冲在各轮间交替使用，形状为 $[b_{\rm cap},\gscan,c]$ 与 $[b_{\rm cap},\gscan/2,c]$；其中 $b_{\rm cap}$ 为预留查询容量：
\begin{equation}
 S_{\rm pairs}=8b_{\rm cap}c\left(\gscan+\frac{\gscan}{2}\right)
 \quad\text{bytes},\qquad \gscan>1.
 \label{eq:5-group-workspace}
\end{equation}
计数另占 $4b_{\rm cap}\gscan$ bytes，已预分配的共享界另占 $8b_{\rm cap}$ bytes，候选和其他工作区另计。这些分配在图捕获前扩容，之后持续保留以复用。归并 block 的共享内存保存列表长度，而非全部键。

短组或空组只贡献有效前缀。排序与输出使用不同补位键：
\[
 K_{\rm sort\ pad}=(+\infty,\mathrm{INT\_MAX}),\qquad
 K_{\rm output\ pad}=(+\infty,-1).
\]
有效位置非负且小于 $\mathrm{INT\_MAX}$；负输出位置表示未使用条目。空组计数为零，$\gscan=1$ 无需归并，直接写出。

单 block 路径 $\gscan=0$ 最多保留 $c$ 个最小计算分数，但截断处完全并列时不固定保留哪些候选。对非空 $V$，令 $h=\min(c,|V|)$，$d_h$ 为其第 $h$ 个分数：
\[
 \{x:d_x<d_h\}\subseteq B\subseteq\{x:d_x\leq d_h\},\qquad |B|=h.
\]
因此，所有严格更小的分数都被保留，保留项也都不大于截断值；并列候选身份可以变化。评测构建关闭了可选的配额式近似选择器。\Cref{sec:7-strict-selection} 通过 CPU 排序 GPU 已计算分数，检查这一距离顺序约束及分组完整键排序。
